\chapter{2008年江苏卷}

\section{填空题}

\begin{problem}
若函数 \( f(x) = \cos \left( \omega x - \frac{\pi}{6} \right) (\omega > 0) \) 最小正周期为 \( \frac{\pi}{3} \)，则 \( \omega = \)  \fillinblank{}.
\end{problem}

\begin{problem}
若将一颗质地均匀的骰子（一种余面上分别标有 1, 2, 3, 4, 5, 6 个点的正方体玩具）先后抛掷 2 次，则出现向上的点数之和为 4 的概率是 \fillinblank{} \fillinblank{}.
\end{problem}

\begin{problem}
若将复数 \(\frac{1 + \i}{a + b - \i}\) 表示为 \(a + b(\i, a, b \in \R, \i\) 是虚数单位）的形式，则 \(a + b = \frac{1}{1 - \i}\) \fillinblank{}.
\end{problem}

\begin{problem}
设集合 \(A = \left\{x \mid (x - 1)^2 < 3x + 7, x \in \R\right\}\)，则集合 \(A \cap \Z\) 中有 个元素 \fillinblank{}.
\end{problem}

\begin{problem}
已知向量 \(a\) 与 \(b\) 的夹角为 \(120^{\circ}\)，\(|a| = 1\)，\(|b| = 3\)，则 \(|5a - b| = \)  \fillinblank{}.
\end{problem}

\begin{problem}
在平面直角坐标系 \(xOy\) 中，设 \(D\) 是横坐标与纵坐标的绝对值均不大于2的点构成的区域，\(E\) 是到原点的距离不大于1的点的构成的区域. 向 \(D\) 中随机投一点，所投的点的落在 \(E\) 中的概率是 \fillinblank{} \fillinblank{}.
\end{problem}

\begin{problem}
某地区为了解 \(70 \sim 80\) 岁老人的日平均睡眠时间（单位：h），随机选择了 50 位老人进行调查. 下表是 50 位老人日睡眠时间频率分布表： 序号 (t) 分组 (睡眠时间) 组中值 (Gt) 频数 (人数) 频率 (Hz) 1 (4,5) 4.5 6 0.12 2 (5,6) 5.5 10 0.20 3 (6,7) 6.5 20 0.40 4 (7,8) 7.5 75 0.20 5 (8,9) 8.5 4 0.08 序号 (t) 分组 (睡眠时间) 组中值 (G) 频数 (人数) 频率 (Hz) 1 (4,5) 4.5 6 0.12 2 (5,6) 5.5 10 0.20 3 (6,7) 6.5 20 0.40 4 (7,8) 7.5 10 0.20 5 (8,9) 8.5 4 0.08

在上述统计数据的分析中，一部分计算见算法流程图，则输出的 S 的值为 \fillinblank{} \fillinblank{}.
\end{problem}

\begin{problem}
设直线 \( y = \frac{1}{2} x + b \) 是曲线 \( y = \ln x (x > 0) \) 的一条切线，则实数 \( b \) 的值为 \fillinblank{} \fillinblank{}.
\end{problem}

\begin{problem}
如图，在平面直角坐标系 \(xOy\) 中，设三角形 \(ABC\) 的顶点分别为 \(A(0, a)\)，\(B(b, 0)\)，\(C(c, 0)\)，点 \(P(0, p)\) 为线段 \(AO\) 上的一点 (异于端点). 这里 \(a, b, c, p\) 为非零实数. 设直线 \(BP, CP\) 分别与边 \(AC, AB\) 交于点 \(E, F\). 某同学已正确求得 \(OE\) 的方程: \(\left(\frac{1}{b} - \frac{1}{c}\right)x + \left(\frac{1}{p} - \frac{1}{a}\right)y = 0\). 请你完成直线 \(OF\) 的方程: \(\fillinblank{}\cdot x + \left(\frac{1}{p} - \frac{1}{a}\right)y = 0\).
\end{problem}

\begin{problem}
将全体正整数排成一个三角形数阵: \[ \begin{array}{r r r r r r r r r} \& \& \& \& \& \& \& \& 1 \\ \& \& \& \& \& \& \& \& 2 \& 3 \\ \& \& \& 4 \& 5 \& 6 \\ \& \& \& 7 \& 8 \& 9 \& 1 0 \\ \& \& \& \& \& \& \& \dots \end{array} \] 根据以上排列规律，数阵中第 \(n\) （\(n \geqslant 3\)）行的从左向右的第3个数是 \fillinblank{} \fillinblank{}.
\end{problem}

\begin{problem}
设 \(x, y, z\) 为正实数，满足 \(x - 2y + 3z = 0\)，则 \(\frac{y^2}{xz}\) 的最小值为 \fillinblank{} \fillinblank{}.
\end{problem}

\begin{problem}
在平面直角坐标系 \(xOy\) 中，设椭圆 \(\frac{x^2}{a^2} + \frac{y^2}{b^2} = 1 (a > b > 0)\) 的焦距为 \(2c\)，以点 \(O\) 为圆心，\(a\) 为半径作圆 \(M\). 若过点 \(\left(\frac{a^2}{c}, 0\right)\) 所作圆 \(M\) 的两条切线互相垂直，则该椭圆的离心率 \(e = \)  \fillinblank{}.
\end{problem}

\begin{problem}
满足条件 \(AB = 2, AC = \sqrt{2} BC\) 的三角形 \(ABC\) 的面积的最大值是 \fillinblank{} \fillinblank{}.
\end{problem}

\begin{problem}
设函数 \( f(x) = ax^2 - 3x + 1 (x \in \R) \)，若对于任意 \( x \in [-1, 1] \)，都有 \( f(x) \geqslant 0 \) 成立，则实数 \( x \) 的值为 \fillinblank{} \fillinblank{}.
\end{problem}

\section{解答题}

\begin{problem}
如图，在平面直角坐标系 \(xOy\) 中，以 \(Ox\) 轴为始边做两个锐角 \(\alpha, \beta\)，它们的终边分别与单位圆交于 \(A, B\) 两点. 已知 \(A, B\) 的横坐标分别为 \(\frac{\sqrt{2}}{10}\)，\(\frac{2\sqrt{5}}{5}\). (1) 求 \(\tan (\alpha + \beta)\) 的值; (2) 求 \(\alpha + 2\beta\) 的值.
\end{problem}

\begin{problem}
如图，在四面体 \(ABCD\) 中，\(CB = CD\)，\(AD \perp BD\)，点 \(E\)、\(F\) 分别是 \(AB\)、\(BD\) 的中点. 求证: (1) 直线 \(EF //\) 平面 \(ACD\); (2) 平面 \(EFC \perp\) 平面 \(BCD\).
\end{problem}

\begin{problem}
如图，某地有三家工厂，分别位于矩形 \(ABCD\) 的两个顶点 \(A\)、\(B\) 及 \(CD\) 的中点 \(P\) 处，已知 \(AB = 20 \mathrm{~km}\)，\(BC = 10 \mathrm{~km}\). 为了处理三家工厂的污水，现要在矩形区域上 (含边界)，且与 \(A\)、\(B\) 等距离的一点 \(O\) 处建造一个污水处理厂，并铺设三条排污管道 \(AO\)、\(BO\)、\(PO\). 记排污管道的总长度为 \(y\) \(\mathrm{km}\). (1) 按下列要求写出函数关系式: \circled{1} 设 \(\angle BAO = \theta\) (md)，将 \(y\) 表示为 \(\theta\) 的函数; \circled{2} 设 \(OP = x(\mathrm{km})\) ，将 \(y\) 表示为 \(x\) 的函数； (2) 请你选用 (1) 中的一个函数关系式，确定污水处理厂的位置，使铺设的排污管的总长度最短.
\end{problem}

\begin{problem}
在平面直角坐标系 \(xOy\) 中，设二次函数 \(f(x) = x^2 + 2x + b (x \in \R)\) 的图象与两坐标轴有三个交点，经过这三个交点的圆记为 \(C\). (1) 求实数 b 的取值范围; (2) 求圆 C 的方程; (3) 问圆 C 是否经过定点 (其坐标与 b 无关)\(\perp\) 请证明你的结论.
\end{problem}

\begin{problem}
(1) 设 \(a_1, a_2, \ldots, a_n\) 是各项均不为零的 \(n (n \geqslant 4)\) 项等差数列，且公差 \(d \neq 0\). 若将此数列删去某一项得到的数列（按原来的顺序）是等比数列. \circled{1} 当 n = 4 时，求 \( \frac{a_{1}}{a_{2}} \) 的数值; \circled{2} 求 n 的所有可能值. (2) 求证: 对于给定的正整数 \(n\) (\(n \geqslant 4\))，存在一个各项及公差都不为零的等差数列 \(b_{1}, b_{2}, \cdots, b_{n}\)，其中任意三项 (按原来顺序) 都不能组成等比数列.
\end{problem}

\begin{problem}
四选二.【A】如图，设 \(\triangle ABC\) 的外接圆的切线 \(AE\) 与 \(BC\) 的延长线交于点 \(E\)，\(\angle BAC\) 的平分线与 \(BC\) 交于点 \(D\). 求证: \(ED^{2} = EC \cdot EB\).

【B】在平面直角坐标系 \(xOy\) 中，设椭圆 \(4x^{2} + y^{2} = 1\) 在矩阵 \(\mathbf{A} = \begin{bmatrix} 2 & 0 \\ 0 & 1 \end{bmatrix}\) 对应的变换下得到曲线 F，求 F 的方程.
\end{problem}

\begin{problem}
如图，设动点 \(P\) 在棱长为 1 的正方体 \(ABCD - A_{1}B_{1}C_{1}D_{1}\) 的对角线 \(BD_{1}\) 上，记 \(\frac{D_{1}P}{D_{1}B} = \lambda\). 当 \(\angle APC\) 为钝角时，求 \(\lambda\) 的取值范围.
\end{problem}

\begin{problem}
已知函数 \( f_{1}(x) = 3^{x - p + 1} \)，\( f_{2}(x) = 2 \cdot 3^{x - p + 1} (x \in \R, p_{1}, p_{2} \) 为常数). 函数 \( f(x) \) 定义为: 对每个给定的实数 \( x, f_{1}(x) = \left\{ \begin{array}{l} f_{1}(x), \text{若 } f_{1}(x) \leqslant f_{2}(x), \\ f_{2}(x), \text{若 } f_{1}(x) > f_{2}(x). \end{array} \right. \) (1) 求 \( f(x) = f_{1}(x) \) 对所有实数 \( x \) 成立的充分必要条件 (用 \( p_{1}, p_{2} \) 表示); (2) 设 \( a, b \) 是两个实数，满足 \( a < b \)，且 \( p_{1}, p_{2} \in (a, b) \). 若 \( f(a) = f(b) \)，求证: \( f(x) \) 在区间 \( [a, b] \) 上的单调增区间的长度之和为 \( \frac{b - a}{2} \). (闭区间 \( [m, n] \) 的长度定义为 \( n - m \))
\end{problem}

\begin{problem}
四选二.【A】如图，设 \(\triangle ABC\) 的外接圆的切线 \(AE\) 与 \(BC\) 的延长线交于点 \(E\)，\(\angle BAC\) 的平分线与 \(BC\) 交于点 \(D\). 求证: \(ED^2 = EC \cdot EB\).
\end{problem}

\begin{problem}
如图，设动点 \(P\) 在棱长为 1 的正方体 \(ABCD - A_{1}B_{1}C_{1}D_{1}\) 的对角线 \(BD_{1}\) 上，记 \(\frac{D_{1}P}{D_{1}B} = \lambda\). 当 \(\angle APC\) 为钝角时，求 \(\lambda\) 的取值范围.

【B】在平面直角坐标系 \(xOy\) 中，设椭圆 \(4x^{2} + y^{2} = 1\) 在矩阵 \(\mathbf{A} = \begin{bmatrix} 2 & 0 \\ 0 & 1 \end{bmatrix}\) 对应的变换下得到曲线 \(F\)，求 \(F\) 的方程. 【C】在平面直角坐标系 \(xOy\) 中，设 \(P(x,y)\) 是椭圆 \(\frac{x^2}{3} + y^2 = 1\) 上的一个动点，求 \(S = x + y\) 的最大值. 【D】设 \(a, b, c\) 为正实数，求证：\(\frac{1}{a^3} + \frac{1}{b^3} + \frac{1}{c^3} + abc \geqslant 2\sqrt{3}\).
\end{problem}

\begin{problem}
请先阅读: 在等式 \(\cos 2x = -2\cos^2 x - 1\) (\(x \in \R\)) 的两边对 \(x\) 求导 \((\cos 2x)^2 = (2\cos^2 x - 1)^2\). 由求导法得到 \((- \sin 2x) \cdot 2 = 4\cos x \cdot (-\sin x)\)，化简得等式 \(\sin 2x = 2\sin x\cos x\). (1) 利用上述想法 (或者其他方法)，试由等式 \((1 + x)^n = C_n^0 + C_n^1 x + C_n^2 x^2 + \cdots + C_n^n x^n\) (\(x \in \R\). 整数 \(n \geqslant 2\)) 证明: \(\eta((1 + x)^n - 1) = \sum_{k=2}^{n} k C_n^k x^{k-1}\); (2) 对于正整数 \(n \geqslant 3\)，求证: \circled{1} \(\sum_{k=1}^{n} (-1)^k k C_n^k = 0\); \circled{2} \(\sum_{k=1}^{n} (-1)^k k^2 C_n^k = 0\); \circled{3} \(\sum_{k=0}^{n} \frac{1}{k+1} C_n^k = \frac{2^{n+1}-1}{n+1}\).
\end{problem}

